% Gerado por tools/tei2latex.py a partir de revisado.xml. Regenerar, não editar.

%%>>> parte-1

%%>>> cap-1

%%>>> art-8

%%>>> t41

\pg{47} \tese{41}{O pensamento não é uma funcção do cerebro.}

a) Toda a funcção do cerebro é sensivel, concreta e determinada pelo tempo e pelo espaço. Ora o pensamento, como dissemos, escapa á percepção dos sentidos, é universal e pre-

scinde de todo o tempo e de todo o espaço.

b) Se o pensamento fosse uma funcção do cerebro, entre a faculdade pensante e o proprio cerebro deveria existir \emph{sempre} uma equação completa e rigorosa. Ora a experiencia demonstra que similhante equação não existe \emph{sempre}.

Logo o pensamento não é uma funcção do cerebro \footnote{Os \emph{materialistas}, para provar que o pensamento é uma funcção do cerebro, recorrem á \emph{anatomia} e á \emph{physiologia}. Na opinião d'elles, a \emph{anatomia} diz que entre o pensamento e o cerebro existe \emph{sempre} uma equação \emph{completa} e \emph{rigorosa}, e a \emph{physiologia} diz que o pensamento é uma \emph{funcção} do cerebro.\par Respondemos que, embora entre o cerebro e o pensamento existisse \emph{sempre} uma rigorosa e perfeita equação, isto só provaria que o cerebro é uma\pgn{48}\emph{condição} do pensamento, e não que é \emph{causa}. Mas é falso que a \emph{anatomia} ensine que entre o cerebro e o pensamento existe aquella equação; como é falso que a \emph{physiologia} ensine que o pensamento é uma funcção do cerebro.\par Em defeza d'estas nossas proposições, provaremos primeiramente que a \emph{anatomia} não demonstra existir entre o cerebro e o pensamento uma \emph{rigorosa} equação, e que a \emph{physiologia} não ensina ser o cerebro o orgão do pensamento; em seguida, provaremos que o cerebro é apenas o orgão das \emph{condições} do pensamento, e finalmente faremos a critica d'algumas asserções dos mais famigerados materialistas.\par O estudo \emph{anatomico} do cerebro não é recente, pois vem de edades remotas; mas o estudo das suas relações com o pensamento é moderno e principiou com a reacção, que tinha por fim recompor o homem e restituir-lhe a unidade pessoal, que \emph{Descartes} tão inconsideradamente tinha quebrado.\par a) \emph{Gall} foi quem primeiro tentou resolver o problema das relações entre o cerebro e o pensamento. Defensor acerrimo da unidade pessoal do homem, o naturalista allemão indicou o caminho, que um sabio deve seguir nas suas investigações ácerca do homem, mas não conseguiu o seu fim. Defendeu a espiritualidade da alma humana, mas inventou um systema, que devia subministrar muitas armas aos materialistas.\par O \emph{phrenologismo} foi recebido com enthusiasmo; pois uma simples observação do craneo permittia a todos o conhecer as faculdades intellectuaes predominantes no individuo. Segundo \emph{Gall}, no cerebro distinguem-se muitas e diversas proeminencias. Cada uma d'ellas é orgão de uma faculdade primitiva e superior. E, como a forma do craneo corresponde perfeitamente á forma do cerebro, legitimamente nas saliencias do craneo podemos reconhecer e distinguir as faculdades da alma. --- \emph{Spurzheim}, discipulo de \emph{Gall}, encontrou mais proeminencias no cerebro, e deu ás faculdades da alma uma disposição ou localisação diversa da do seu mestre.\par Esta theoria não tinha uma base solida. \emph{Gall} e \emph{Spurzheim} dividiram o cerebro e localisaram as faculdades a capricho e sem criterio; por isso o \emph{phrenologismo} encontrou opposição entre os sabios, e foi condemnado pela razão, pela sciencia e pela experiencia. --- Foi condemnado pela razão; porque faculdades espirituaes, como são a intelligencia e a vontade, não podem residir em orgãos materiaes. --- Foi condemnado pela sciencia; porque a anatomia demonstra que a conformação exterior do craneo não corresponde á sua conformação interior e por isso á forma do cerebro. --- Foi condemnado pela experiencia; porque no craneo de \emph{Laplace} foi encontrado o orgão da estupidez, no craneo do famoso assassino \emph{Fiesch} foi encontrado o orgão da benevolencia e no do celebre ladrão \emph{Lacenaire} o orgão da liberalidade.\par b) Os materialistas, não encontrando no craneo dos vivos razões para a defeza das suas doutrinas, foram procural-as no craneo dos mortos. Mediram a capacidade do craneo dos mortos; porque d'este modo, diziam, temos a quantidade exacta da massa cerebral, e por isso a quantidade da intelligencia; pois grande volume do encephalo denota necessariamente grande\pgn{49}talento. --- \emph{Büchner} recorre á auctoridade dos chapelleiros, e diz que a quotidiana experiencia d'estes officiaes ensina que a classe culta precisa de chapeus maiores que a classe ignorante.\par Se fosse verdade que uma capacidade maior do craneo denota sempre maior intelligencia, os craneos dos antepassados deviam ter uma capacidade inferior á dos nossos; pois os materialistas dizem que os antepassados eram dotados de uma intelligencia mediocre, muito inferior á dos homens da nossa epoca. Todavia a experiencia mostra o contrario. Encontraram-se \emph{seis} craneos, nas grutas de Lozère, na França, pertencentes a uma epoca remotissima, a que deram o nome de \emph{epoca da pedra polida}. Pois bem, cada um d'esses craneos tinha uma capacidade media de 1606 centimetros cubicos; ao passo que da experiencia, feita em 67 craneos de individuos parisienses, do ultimo seculo, resultou que estes craneos tinham uma capacidade media de 1559c. c., isto é, 47c. c. menos que os craneos da \emph{pedra polida}.\par Mas façamos uma comparação entre os craneos dos individuos de diversas nações na epoca presente. A capacidade do craneo dos Esquimós é de 1546c.c., a dos Neo-Zelandezes é de 1497c.c., a dos Japonezes de 1486c.c., a dos Europeus de 1467c.c. Se fosse verdadeira a opinião dos materialistas, os Esquimós seriam os mais intelligentes do mundo, e os Europeus só occupariam o quarto grau na escala da intelligencia; o que é falso.\par c) \emph{Topinard} e os seus discipulos ficaram deveras embaraçados á vista d'este resultado, como de um phenomeno imprevisto; mas não se deram por vencidos. Mudaram apenas o campo da contenda, e disseram que não é a capacidade do craneo que se deve medir, mas é o proprio cerebro que se deve pesar; pois o peso do cerebro augmenta com o grau do desenvolvimento intellectual; quanto mais desenvolvido fôr o cerebro, tanto mais poderosa será a intelligencia.\par Este recurso não salva o materialismo. Appellemos para a \emph{anatomia}, que nos dá os seguintes resultados:\par O cerebro de um elephante pesa em media . . 3000 gr.\par O do delphim. . . . . . . . . . . . . . 1800 »\par O da baleia. . . . . . . . . . . . . . . 1500 »\par O do homem. . . . . . . . . . . . . . 1300 «\par Se, pois, o desenvolvimento do cerebro fosse proporcionado ao desenvolvimento intellectual, o elephante, o delphim e a baleia seriam mais intelligentes do que o homem.\par d) Os materialistas dizem que não se deve attender ao peso \emph{absoluto} do cerebro, mas sim ao peso \emph{relativo}.\par Respondemos que este recurso é perfeitamente inutil e vão. Prescindindo dos outros animaes, que formam a escala zoologica, e considerando só os mammiferos, temos os seguintes resultados. No cavallo, o peso do cerebro está para o peso do corpo como 1 para 400; no macaco, como 1 para 100; no homem, como 1 para 47; no saguim, como 1 para 28; no canario, como 1 para 14. Se, pois, valesse o argumento dos materialistas, o saguim\pgn{50}e o canario deveriam ser mais intelligentes do que o homem, pois o peso do cerebro d'aquelles animaes é \emph{relativamente} maior que o do homem.\par e) Os materialistas, porém, querem limitar-se ao peso \emph{relativo} do cerebro das raças humanas, e dizem que, na mesma especie humana, o peso do cerebro das diversas raças está em rigorosa proporção com os differentes graus de civilisação e por isso de intelligencia.\par Similhante asserção, que, ha pouco tempo, era tida como um axioma scientifico, está hoje abandonada, como contraria á experiencia.\par Eis o peso do cerebro de alguns povos.\par O cerebro dos Escocezes pesa em media. . . . 1417 gr.\par O dos Insulares das Carolinas. . . . . . . . 1402 »\par O dos Esquimós. . . . . . . . . . . . . . 1398 «\par O dos Inglezes. . . . . . . . . . . . . . 1388 »\par O dos Bavaros. . . . . . . . . . . . . . 1375 »\par O dos Francezes. . . . . . . . . . . . . 1359 »\par O dos Annamitas. . . . . . . . . . . . . 1341 »\par O dos Italianos. . . . . . . . . . . . . . 1308 »\par Na opinião, pois, dos materialistas, os Insulares das Carolinas, os Esquimós e os Annamitas deveriam ser mais intelligentes do que os Inglezes, os Bavaros, os Francezes e os Italianos.\par f) A falsidade do systema materialista torna-se ainda mais evidente, se considerarmos o cerebro no seu desenvolvimento. --- Eis os resultados anatomicos, segundo as experiencias de \emph{Boyd}.\par No nascimento, o cerebro humano pesa. . . . 331 gr.\par Na edade de 3 mezes . . . . . . . . . . . . 493 »\par De 3 a 6 mezes . . . . . . . . . . . . . . 603 »\par De 6 mezes a 1 anno. . . . . . . . . . . . 777 »\par De 1 a 2 annos. . . . . . . . . . . . . . 942 «\par De 2 a 4 annos. . . . . . . . . . . . . . 1097 »\par De 4 a 7 annos. . . . . . . . . . . . . . 1140 »\par De 7 a 14 annos. . . . . . . . . . . . . 1302 »\par De 14 a 20 annos. . . . . . . . . . . . . 1324 »\par De 20 a 30 annos. . . . . . . . . . . . . 1357 »\par De 30 a 40 annos. . . . . . . . . . . . . 1366 »\par De 40 a 50 annos. . . . . . . . . . . . . 1352 »\par De 50 a 60 annos. . . . . . . . . . . . . 1343 »\par De 60 a 70 annos. . . . . . . . . . . . . 1315 »\par De 70 a 80 annos. . . . . . . . . . . . . 1289 »\par De 80 a 90 annos. . . . . . . . . . . . . 1284 »\par Conhecidos estes resultados, tiramos as seguintes conclusões.\par 1.) Se o cerebro fosse o orgão do pensamento, a criança, que tem um anno de edade, já deveria pensar, porque tem 777 grammas de cerebro, isto é, um volume superior ao volume do macaco já adulto, que, segundo as doutrinas dos materialistas, é dotado de intelligencia.\par 2.) De mais, se o grau da intelligencia estivesse na razão directa do peso\pgn{51}do cerebro, todas as crianças de 12 annos, cujo cerebro pesa 1260 grammas, deveriam ser mais intelligentes do que o celebre \emph{Gambetta}, cujo cerebro apenas pesava 1160 grammas, isto é, 140 gr. menos que a media.\par 3.) Finalmente, se do peso do cerebro dependesse o grau da intelligencia, então o homem, depois dos 40 annos, quando o peso do cerebro principia a declinar, deveria declinar nas faculdades intellectuaes. Mas não é assim; porque todos sabem que, geralmente, depois dos 40 annos principia a edade da reflexão, do juizo e do conselho.\par Verdade é esta tão evidente, que um dos mestres materialistas, \emph{Broca}, apertado pelas objecções de \emph{Gratiolet}, chegou a fazer a seguinte declaração, que está em contradicção aberta com as doutrinas materialistas: » Ninguem pretendeu, nem pode pretender, que exista uma relação \emph{absoluta} entre o desenvolvimento da intelligencia e o volume do cerebro. Protestei sempre contra este absurdo. O sabio nem sequer póde pensar em que, medindo o cerebro, meça a intelligencia ».\par Todos estes desmentidos da \emph{anatomia} deviam ser mais que sufficientes para obstar a que os materialistas apregoassem as suas doutrinas em nome da sciencia. Mas não succedeu assim. Obrigados a reconhecer que nem nas proeminencias e na cavidade do craneo, nem no peso \emph{absoluto} ou \emph{relativo} do cerebro encontravam uma base para o seu edificio, os materialistas disseram que, em fim, o volume do cerebro é um dos factores da intelligencia, mas não é \emph{unico}; pois o grau mais ou menos elevado da intelligencia depende da symetria dos hemispherios cerebraes, da quantidade da substancia cinzenta e de phosphoro, e da grandeza dos lóbos frontaes.\par Examinemos estas affirmações.\par g) Foi o naturalista \emph{Bichat}, quem fez depender o maior ou menor desenvolvimento da intelligencia da symetria, mais ou menos perfeita, dos hemispherios cerebraes.\par Mas esta opinião morreu com o proprio auctor. Na autopsia feita ao cadaver de \emph{Bichat}, um dos hemispherios cerebraes foi encontrado atrophiado; e todavia esse anatomista era intelligentissimo. Em seguida, a sciencia demonstrou que o atrophiamento parcial ou total de um só hemispherio póde deixar e muitas vezes deixa intacta a intelligencia.\par h) Appellando para a substancia \emph{cinzenta}, os materialistas disseram que os que têm mais substancia \emph{cinzenta}, têm mais intelligencia. E, como a substancia \emph{cinzenta} forma a camada cortical do cerebro, sustentaram que o grau da intelligencia será mais ou menos elevado, segundo forem mais ou menos salientes as \emph{circumvoluções}.\par Não gastaremos muitas palavras para refutar esta doutrina. Só diremos que, se devessemos acreditar nas palavras dos materialistas, então mais intelligentes do que todos os outros animaes, mais intelligentes do que o proprio homem seriam o carneiro e o... jumento; porque as \emph{circumvoluções} do cerebro d'estes dois animaes são salientissimas, e, pelo contrario, o mais estupido de todos os animaes seria o castor, que tem o cerebro liso.\par i) Outros materialistas, como \emph{Moleschott}, \emph{Couerbe}, \emph{Feuerbach}, attribuem o maior ou menor desenvolvimento da intelligencia á maior ou menor quantidade\pgn{52}de \emph{phosphoro} contido no cerebro; e dizem que a abundancia de \emph{phosphoro} no cerebro produz genio; o seu excesso, loucura; e a sua ausencia ou escassez, idiotismo.\par Esta opinião foi refutada por muitos sabios, e especialmente por \emph{Frémy}. As analyses chimicas, realizadas no cerebro de muitos loucos, mostraram que elle não tem mais \emph{phosphoro} do que o dos outros homens. --- De mais, o cerebro dos peixes e o do animal suino têm mais phosphoro, que o do homem; e, todavia, ninguem dirá que esses animaes são mais intelligentes do que o homem. --- Finalmente, o phosphoro desenvolve-se no cerebro dos animaes depois da sua morte, e por isso não póde ser a causa da intelligencia.\par j) Outros materialistas dizem que a intelligencia está localisada nos \emph{lóbos frontaes}. --- \emph{Topinard} diz: « Nos lóbos frontaes residem as mais altas faculdades. Lóbos volumosos, e por isso fronte espaçosa, revelam grande intelligencia; lóbos atrophiados, e por isso fronte baixa, são indicio de uma intelligencia mediocre ou nulla ».\par Respondemos que o desenvolvimento dos lóbos frontaes não é indicio seguro do desenvolvimento da intelligencia. Se em alguns idiotas esses lóbos estavam atrophiados, em outros conservavam o seu estado normal. Se Walter \emph{Scott} e \emph{Cuvier} possuiam uma fronte alta, o celebre \emph{Lacépède} e o General \emph{Lamarque} eram dotados de uma fronte baixa.\par k) Outros materialistas, como \emph{Newmann de Leipzig}, collocaram a intelligencia nos lóbos occipitaes, porque no cerebro de muitos loucos esses lóbos estavam mais atrophiados do que os frontaes.\par Esta opinião, como a antecedente, está desmentida pela moderna theoria das \emph{localisações}, --- theoria que tem por auctores dois physiologistas allemães, \emph{Fritsch} e \emph{Hitzig}. Estes sabios demonstraram até á evidencia que os \emph{centros motores} residem quasi exclusivamente no lóbo parietal e na metade posterior do lóbo frontal, e que a sensibilidade reside na parte posterior dos hemispherios e especialmente nos lóbos occipital e temporo-esphenoidal. D'este modo todo o cerebro só é orgão do movimento e da sensibilidade. A intelligencia, pois, não tem orgão, e por isso não é uma funcção do cerebro. --- E' verdade que os materialistas dão aos \emph{centros motores} o nome de \emph{centros psychico-motores}, e fazem d'esses centros o orgão da intelligencia; mas esta pretenção foi eliminada pela sciencia, que nas camadas corticaes reconheceu o unico e exclusivo officio de \emph{centros motores}.\par A \emph{anatomia}, pois, longe de favorecer o materialismo, refuta-o.\par Concluiremos com Waitz; « Duvidamos que haja um só individuo, dotado de alguns conhecimentos physiologicos, que se persuada de que o espirito se avalie pelo peso do cerebro, se meça pela grandeza do craneo, pelo desenvolvimento comparativo das differentes partes da massa encephalica, e de que as innumeraveis aptidões ou disposições intellectuaes, moraes e sentimentaes se desenhem nas protuberancias da superficie do craneo ».\par Os materialistas recorrem tambem á \emph{physiologia}, e, em nome d'esta\pgn{53}sciencia, proclamam que o cerebro é orgão do pensamento, pois que está averiguado que o cerebro \emph{trabalha}, quando o homem pensa.\par Vejamos se verdadeiramente a \emph{physiologia} dá razão aos materialistas.\par E' certo que o cerebro \emph{trabalha}, quando o homem pensa.\par Agora pergunta-se: O cerebro será \emph{o orgão}, a \emph{causa efficiente} do pensamento, ou apenas uma \emph{condição} ?\par Os materialistas dizem que o cerebro é o \emph{orgão}, a \emph{causa efficiente} do pensamento; nós, pelo contrario, sustentamos que é apenas uma \emph{condição}.\par Será verdade que a \emph{physiologia} ensina ser o cerebro o \emph{orgão} do pensamento? --- Respondemos que não. Temos em nosso favor o testemunho dos celebres \emph{Gavarret} e \emph{Griesinger}, e dos insuspeitos \emph{Du Bois - Reymond} e \emph{Tyndall}.\par \emph{a}) \emph{Gavarret} assim se exprime: « Da parte do cerebro ha um augmento de actividade e de combustão, e uma producção de calor; e, ao mesmo tempo, dá-se uma manifestação psychica. Entre aquelle trabalho interior e essa manifestação existe uma coincidencia perfeita; o primeiro é evidentemente uma \emph{condição} da segunda. Mas haverá, além d'esta, uma outra relação entre os dois phenomenos? Existirá uma medida commum para a quantidade de calor que se exhala e um pensamento que se concebe? Emquanto se não encontrar essa medida commum, não temos o direito de affirmar --- \emph{que o trabalho cerebral e a concomitante manifestação psychica só differem na forma}, --- \emph{que estes dois productos têm a mesma natureza}, --- \emph{que o primeiro é causa sufficiente do segundo} « (\emph{Les phénomènes physiques de la vie}, 1869).\par \emph{b}) \emph{Griesinger} diz; « Como é que um phenomeno material e physico, que se realiza nas fibras nervosas ou nas cellulas ganglionarias, póde tornar-se uma idea, ou um acto de consciencia? Isto é absolutamente incomprehensivel » (\emph{Traité des maladies mentales}, trad. Doumic. 1865).\par \emph{c}) \emph{Du Bois-Reymond} escreve: « Ainda que conhecessemos intrinsecamente o cerebro, os phenomenos intellectuaes seriam sempre incomprehensiveis.... O mais profundo conhecimento do encephalo só nos manifesta a materia em movimento. Mas nenhuma disposição, nenhum movimento de partes materiaes póde servir-nos de ponte de passagem para o dominio da intelligencia. . . Á primeira vista e por uma observação superficial parece que, pelo conhecimento dos factos materiaes no cerebro, poderiamos comprehender certos factos de ordem espiritual...; mas a mais breve reflexão mostra que isto é uma illusão. Apenas podemos obter noções sobre as circumstancias da vida espiritual, que se tornam manifestas pelas impressões dos sentidos, mas não sobre a natureza da vida espiritual por meio d'essas circumstancias... De modo nenhum se póde comprehender como a consciencia possa nascer da combinação de certo numero de atomos de carbone, hydrogenio, azote, oxygenio, etc. Não é preciso demonstrar a impossibilidade actual e permanente de comprehender os elevados phenomenos espirituaes pela supposta mechanica dos atomos cerebraes... Os nossos conhecimentos naturaes estão encerrados entre dois limites: de um lado a incapacidade para comprehender a materia e a força, do outro a impossibilidade de deduzir\pgn{54}das condições materiaes os phenomenos espirituaes: estes limites estão determinados irrevogavelmente » (\emph{Sur les bornes de la philos. naturelle}).\par \emph{d}) \emph{Tyndall} diz: « E' impossivel conceber a passagem da physica do cerebro ao facto correspondente da consciencia intima das sensações, dos pensamentos, das emoções. Se nos concedessem que um pensamento determinado e uma determinada acção do cerebro são factos simultaneos, não só não teriamos demonstrado a existencia do orgão intellectual, mas nem sequer a existencia de um rudimento visivel d'este orgão. Ainda quando a nossa intelligencia e os nossos sentidos fossem bastante aperfeiçoados, vigorosos e illuminados para deixarem ver e sentir as ultimas particulas do cerebro, e seguil-as nos seus movimentos; ainda quando tivessemos a consciencia dos estados correspondentes do pensamento e do sentimento, estariamos sempre muito longe da solução d'este problema. Qual é a relação entre as operações physicas e os factos da consciencia? O abysmo entre as duas classes de phenomenos é e será sempre insuperavel » (\emph{Les mondes}).\par Do mesmo modo se exprimem \emph{Jaccoud} (\emph{Traité de pathologie interne}), Potain (Art. \emph{Cerveau}, no \emph{Dict. encycl. des sciences med.}), Ferrier (\emph{Fonctions du cerveau}), Vierordt (\emph{Le physiologie de l'homme}).\par Estas declarações tiram ao materialismo toda a apparencia scientifica. A physiologia constata o facto, diz que existe uma relação entre o trabalho do cerebro e o pensamento, mas não póde affirmar se essa relação é \emph{causal} ou apenas \emph{condicional}. --- Para a physiologia, o pensamento é incomprehensivel- Na verdade, esta sciencia, para attingir o seu objecto, emprega instrumentos materiaes, --- o escalpello, o microscopio, a balança; mas o pensamento não tem partes, em que se possa dividir, nem côr, que se veja, nem quantidade, que se pese. Quem viu as fracções da idea de \emph{justiça}, ou quem póde dizer se o pensamento é branco ou vermelho, se pesa dois ou tres grammas?\par Nem appellem os materialistas para as experiencias de \emph{Flourens}. Este sabio constatou simplesmente que a integridade de certas partes do cerebro é indispensavel para o regular exercicio das faculdades intellectuaes. Mas, do facto de uma coisa ser indispensavel para a producção de um phenomeno, poderemos nós concluir que essa coisa é a \emph{causa} ou o orgão do phenomeno ? Não; a sciencia não acceita tal conclusão. Tambem os \emph{tuberculos quadrigemeos} são indispensaveis para o acto regular da visão, e todavia nenhum physiologista póde dizer que elles são o orgão central da visão. Do mesmo modo, se a \emph{physiologia} ensina que o cerebro é indispensavel para nós pensarmos, d'ahi não podemos concluir que o cerebro é o \emph{orgão} do pensamento.\par Concluimos, pois, que nem a \emph{anatomia} nem a \emph{physiologia} podem demonstrar que o cerebro é o \emph{orgão} do pensamento. O pensamento não é uma secreção, nem um movimento, nem outra qualquer funcção do cerebro, mas é uma coisa immaterial, que escapa ao alcance da experiencia sensivel; e por isso, assim como só póde derivar de uma faculdade \emph{espiritual}, assim tambem só póde ser percebido por uma faculdade \emph{espiritual}.\par Se o cerebro não é o \emph{orgão do pensamento}, é, todavia, uma sua \emph{condição}\pgn{55}\emph{indispensavel}. \emph{Lelut}, seguindo as tradições da philosophia escholastica e christã, diz: O cerebro é o orgão das \emph{condições} do pensamento.\par Demonstraremos brevemente a verdade d'esta affirmação.\par O homem não é nem Anjo nem animal, mas participa de um e de outro, é um composto de materia e de espirito, é o \emph{animal racional}.\par Da essencia deriva a operação; esta é effeito d'aquella. Por isso entre uma e outra deve existir uma proporção, de modo que a operação reflicta os caracteres da essencia. Ora, a essencia do homem é composta de dois elementos, --- o \emph{material} e o \emph{espiritual}. Logo o pensamento deve manifestar esses dois elementos, --- o elemento \emph{material} e o elemento \emph{espiritual}.\par No pensamento devemos distinguir duas coisas: o \emph{principio}, de que elle deriva e que é a faculdade pensante, e o \emph{termo}, a que chega e que é o \emph{objecto} pensado. O pensamento não manifesta o \emph{elemento material} no seu \emph{principio}, porque o \emph{principio} do pensamento é a intelligencia, e esta é espiritual. Logo manifesta o elemento \emph{material} no seu \emph{termo}, emquanto o \emph{objecto} é apresentado por uma faculdade material e de um modo concreto.\par Com effeito, o homem, quando nasce, não possue idea alguma. A nossa intelligencia é uma \emph{toboa rasa}, em que tudo se póde escrever, mas nada está escripto. Só depois de alguns annos a intelligencia desenvolve-se, e os primeiros objectos, que conhece, são os corporeos, os materiaes. Mas esses objectos, para que sejam percebidos, devem unir-se á faculdade cognitiva. Como são materiaes, não podem unir-se por si mesmos, mas por uma sua imagem ou representação. Essa imagem, porém, é \emph{concreta} e \emph{singular}, porque representa a essencia do objecto, envolvida nas circumstancias ou condições individuaes de \emph{tempo}, \emph{logar}, \emph{figura}, \emph{côr}, \emph{grandeza}, etc.; e por isso só póde residir na \emph{imaginação}, que é uma \emph{faculdade sensitiva}, visto que na intelligencia só reside o \emph{abstracto}, o \emph{universal}.\par Todavia, a imagem \emph{concreta} e \emph{material}, se não existe na intelligencia, póde comtudo unir-se a esta faculdade. Porquanto, representando aquella imagem a \emph{essencia} do objecto envolvida nas condições materiaes, podemos, por meio da abstracção, prescindir das condições materiaes e só considerar a essencia. A essencia, abstrahindo das condições materiaes, é \emph{universal} e \emph{abstracta}, e, como tal, póde unir-se e une-se á intelligencia.\par Para o conhecimento intellectual, pois, é necessario o exercicio de duas faculdades --- da \emph{imaginação} e da \emph{intelligencia}. O exercicio da \emph{imaginação} é necessario para o da \emph{intelligencia}; porque a intelligencia só póde conhecer os objectos que lhe apresenta a imaginação. Um cego de nascimento comparará a côr vermelha ao som de uma flauta, não porque lhe falta a intelligencia, mas porque na sua imaginação não existe a imagem sensivel da côr vermelha, que elle nunca viu. --- O exercicio da imaginação não só precede, mas acompanha sempre o exercicio da intelligencia. A nossa intelligencia não pode pensar num objecto, sem que a imaginação forme uma imagem sensivel d'esse objecto (ao menos a imagem da \emph{palavra}), na qual de qualquer modo resplandeça o ideal, que contemplamos. Isto verifica-se não só quando o objecto da intelligencia é material, mas tambem quando é immaterial. Não\pgn{56}podemos elevar-nos a \textsc{Deus}, senão por meio do \emph{visivel}; não podemos perceber o \emph{absoluto}, o \emph{infinito}, o \emph{necessario}, senão por meio do \emph{relativo}, do \emph{finito}, do \emph{contingente}.\par Ora a \emph{intelligencia} é uma faculdade \emph{inorganica} e reside na alma; a \emph{imaginação} é uma faculdade \emph{organica} e reside no cerebro. Logo, na operação intellectual, devem concorrer o cerebro com as suas vibrações e movimentos, e a intelligencia com as diversas formas de pensamento. Assim como, sem o exercicio da intelligencia, não existiria o \emph{principio} cognitivo; assim tambem, sem o exercicio do cerebro, faltaria o \emph{objecto}, que deve ser conhecido. O pensamento, pois, deriva \emph{exclusivamente} da intelligencia, mas exige, como \emph{condição indispensavel}, o concurso do cerebro.\par S. Thomaz tinha dito tudo n'estas poucas palavras: « Corpus requiritur ad actionem intellectus, \emph{non sicut organum}, quo talis actio exerceatur, \emph{sed ratione objecti}. Phantasma enim comparatur ad intellectum sicut color ad visum. Sic autem indigere corpore non removet intellectum esse subsistentem. Alioquin animal non esset aliquid subsistens, cum indigeat exterioribus sensibilibus ad sentiendum » (S. T. 1 p. q. 75, a. 2 ad 3).\par Dos principios expostos resultam legitimamente os seguintes corollarios, muito importantes não só para a psychologia mas tambem para a physiologia, --- corollarios, de que os materialistas erradamente se servem para combater a espiritualidade do pensamento e da alma humana.\par 1.º) Se o cerebro não estiver sufficientemente desenvolvido, pelo menos naquellas partes onde se acham localisados os sentidos internos, que concorrem para o acto intellectual e que são especialmente a \emph{memoria} e a \emph{imaginação}, não estará convenientemente desenvolvida a intelligencia.\par 2.º) Enfraquecidas as funcções do cerebro, como nas doenças e em muitos velhos, --- ou alteradas, como na loucura e na embriaguez, --- ou suspensas, como no somno profundo e no deliquio, ficará tambem enfraquecido, alterado ou suspenso o pensamento; porque, não funccionando, ou não funccionando regularmente a sensibilidade interna, a intelligencia ou não terá presente o objecto, ou não o terá presente como elle é na sua realidade. Por isso a loucura, e qualquer outra anomalia intellectual, se não póde dar-se na alma \emph{pianista} de \emph{Descartes}, nem na alma \emph{monada} de \emph{Leibnitz}, completamente isoladas do organismo e cheias de ideas innatas, bem pode dar-se na verdadeira alma humana, que, para adquirir os seus conhecimentos, precisa do concurso das faculdades sensitivas.\par 3.º) Uma lesão, mais ou menos grave, naquellas partes do cerebro, onde se acham localisados os sentidos internos necessarios para o exercicio do pensamento, produzirà tambem uma desordem, mais ou menos grave, na vida intellectual. Este corollario explica sufficientemente certos factos, que a physiologia verifica mas não sabe explicar, como é que a mais pequena lesão do cerebro produz, ás vezes, graves perturbações, ao passo que, outras vezes, graves lesões não produzem o menor desarranjo intellectual.\par 4.º) Quanto mais intenso e prolongado fôr o pensamento, tanto maior será a actividade do cerebro, tanto mais elevado o grau da sua temperatura, tanto mais rapida a circulação do sangue, e tanto mais abundante a combustão da substancia cerebral. Com effeito, attesta-nos a consciencia que a\pgn{57}nossa alma, sem criar novas forças, póde fazer passar da potencia para o acto as que existem em nossos orgãos, e dirigil-as no seu desenvolvimento, provocando movimentos no organismo, por muitos modos e motivos. Ora, se a vontade actua nos musculos e nos orgãos locomotores, a intelligencia actua nos orgãos da memoria e da imaginação, para metter em movimento as imagens sensiveis, que são os materiaes das nossas abstracções, das ideas e de todos os actos intellectuaes. Mas a intelligencia não activa as operações sensitivas dos orgãos cerebraes, sem provocar um augmento nas operações nutritivas d'esses orgãos e por isso sem que o sangue afflua ao cerebro com as riquezas nutritivas, e accelere as reacções chimicas da nutrição.\par 5.º) Num trabalho intellectual muito intenso, sentimos cançaço no cerebro, e especialmente nos \emph{lóbos frontaes}. Dá-se tal phenomeno, não já porque as ideas estejam localisadas nalguma parte do cerebro (as ideas, não tendo figura nem dimensão, não podem residir num sujeito extenso), mas porque ha nesse trabalho phenomenos de ordem sensivel, que são as representações da \emph{imaginação}, localisadas na camada cerebral e provavelmente nos lóbos frontaes.\par Examinemos algumas affirmações e difficuldades dos materialistas.\par a) \emph{Moleschott} escreve : « Se analyso, \emph{pelos processos da chimica}, o cerebro humano, só encontro elementos materiaes e não o espirito; com a chimica decomponho \emph{tudo} em ammoniaco, agua e saes. Logo o principio do pensamento é materia e não o espirito. »\par Concedemos de bom grado que \emph{Moleschott}, \emph{pelos processos da chimica}, só encontra elementos materiaes e não o espirito; pois o espirito não é composto de elementos materiaes. Se um naturalista podesse um dia encontrar na ponta do seu escalpello o espirito, já o espirito não seria espirito. Para se encontrar o espirito, devem empregar-se outros processos que não são os da chimica, devem empregar-se os processos da razão e então esta nos dirá que, realizando-se no homem phenomenos espirituaes, deve tambem existir nelle um principio proporcionado, que é o espirito. --- Concedemos tambem a \emph{Moleschott} que possa decompor tudo o que fôr \emph{chimicamente} composto; mas negamos que possa decompor o que não é nem chimica nem physicamente composto, como é a alma humana.\par b) \emph{Moleschott} continua: « Uma garrafa, contendo carbonato, ammoniaco, chloreto de potassio, phosphato de sodio, cal, etc., é \emph{d'um modo ideal} o \emph{principio vital completo das plantas}, \emph{dos animaes e do proprio homem}. »\par Não é necessario que insistamos na refutação d'esses desvarios. A razão e a experiencia, como provámos, ensinam que o \emph{principio vital} não é materia, nem se identifica com as forças physicas e chimicas da materia. --- \emph{Cl. Bernard} diz: « \emph{Poderia facilmente demonstrar que}, \emph{em physiologia}, o \emph{materialismo a nada guia e nada explica} » (\emph{La science experimentale}, pag. 361).\par c) \emph{Herzen} (\emph{Le cerveau et l'activité cérébrale}) e \emph{Büchner} (\emph{Force at matière}) fazem as seguintes objecções: « A physiologia ensina --- que as operações intellectuaes, como a percepção, o juizo e o raciocinio, demandam tempo, --- que a modificação nas operações organicas importa uma modificação na consciencia e no pensamento, --- e que a razão enfraquece com o enfraquecimento\pgn{58}dos orgãos. Ora estes factos provam que as operações intellectuaes derivam de uma faculdade organica ».\par Repetimos que é a alma espiritual, e só ella, que pensa; mas que, para pensar, precisa do concurso das faculdades sensitivas ou organicas. Por isso no pensamento reflecte-se naturalmente toda e qualquer perturbação ou modificação da percepção sensitiva.\par Este principio explica os factos apontados.\par A percepção, o juizo e o raciocinio demandam tempo, não porque são operações materiaes e organicas, mas porque suppõem uma percepção sensitiva, que é impossivel sem uma operação organica e por isso sem uma mudança material \emph{successiva}, isto é, sem tempo.\par A modificação na consciencia e no pensamento, que deriva da modificação nas operações organicas, explica-se tambem pelo facto dos sentidos fornecerem á intelligencia a materia do pensamento.\par A intelligencia enfraquece com o organismo, porque os orgãos enfraquecidos não fornecem a materia sufficiente para os actos intellectuaes.\par Todavia advertimos que a dependencia entre a faculdade intellectual e o cerebro não é tão absoluta, como se pensa. \emph{Longet} escreve: « Esta dependencia tão natural não é de tal modo absoluta, que se não dêm bastantes exemplos em sentido contrario. Encontram-se meninos debeis que espantam pela precocidade da razão e amplitude do seu espirito, velhos caducos e proximos do sepulchro que conservam intacto o seu juizo, a memoria, a força do genio e o ardor da coragem... Algumas vezes a loucura é acompanhada por uma lesão notavel dos centros nervosos; mas que deverá dizer-se dos casos, em que \emph{Esquirol}, \emph{Leuret} e outros auctores conscienciosos affirmam não ter encontrado vestigio de alteração no cerebro? Os annaes da sciencia fornecem um grande numero de factos perfeitamente observados de alterações profundas na massa cerebral, sem que durante a vida se haja manifestado a mais leve perturbação da intelligencia; têm-se visto porções de cerebro arrancadas, balas que atravessaram esse orgão de parte a parte sem desarranjo do espirito, ao passo que, ás vezes, leves transvasações de sangue, num ponto restricto, promovem a febre, excitam o delirio e produzem rapidamente a morte ». (\emph{Traité de physiologie}, 1873, t. III, pag. 593). --- E basta \emph{uma só excepção} nessa dependencia entre a faculdade intellectual e o cerebro, para se demonstrar que o pensamento não deriva do cerebro; porque, se derivasse realmente, essa excepção não poderia ter logar.\par d) O dr. \emph{Luys}, na sua obra --- \emph{Le cerveau et ses fonctions}, diz que a psychologia não tem razão de ser, e que a physiologia e a anatomia explicam \emph{todo} o homem; pois que o estudo da composição do cerebro pertence á anatomia e o estudo das funcções d'esse orgão é objecto da physiologia.\par Concordamos com o dr. \emph{Luys} em que o estudo das funcções do cerebro é objecto da physiologia. Mas isto nada prova contra a nossa these, se o adversario não demonstrar que o pensamento é uma funcção do cerebro. E estamos convencidos de que nunca o demonstrará. E de que modo o poderia fazer? --- Pela razão? não: porque a razão diz que um phenomeno espiritual, como é o pensamento, não póde ser funcção de um orgão material, como é\pgn{59}o cerebro. --- Pela experiencia? tambem não; porque nenhum physiologista viu funccionar o cerebro, pois que, para isso, seria mister arrancar as meninges, mas, se arrancarmos as meninges, o cerebro deixará de funccionar.\par e) Os materialistas objectam: Perturbado o cerebro, não ha memoria dos conhecimentos passados. Logo não é a alma espiritual, mas é a materia que recorda os pensamentos anteriores.\par Respondemos que, perturbado o cerebro, não ha memoria dos pensamentos passados, não porque seja a materia que pense e recorde, mas porque faltam as imagens sensiveis, e, faltando estas, falta á memoria o seu objecto. --- Pelo contrario, o phenomeno da memoria prova contra os adversarios. \emph{Claudio Bernard} diz: « O corpo humano é um composto de materias, que se renovam incessantemente. Todas as partes do corpo estão sujeitas a um perpetuo movimento de transformação. Cada dia perdeis alguma coisa do vosso ser physico, e pela alimentação substituis o que perdeis. De modo que, num espaço de oito annos, pouco mais ou menos, a vossa carne e os vossos ossos são substituidos por nova carne e por novos ossos, em consequencia das transformações successivas. A mão, com que escreveis hoje, não é inteiramente composta das mesmas moleculas que ha oito annos. A forma é a mesma, mas é uma nova substancia que a completa. O que digo da mão, dil-o-hei do cerebro. A abobada craneana não é occupada pela mesma materia que ha oito annos. Sendo assim, visto que tudo se muda no vosso cerebro, em oito annos, como é que vos lembraes de coisas que foram vistas ou ouvidas ha mais de oito annos? Se as coisas são, como pretendem certos physiologistas, alojadas, incrustadas nos lóbos do vosso cerebro, como é que sobrevivem á desapparição dos lóbos? Esses lóbos não são os mesmos que ha oito annos, e todavia a memoria guardou intacto o seu deposito. E' que no homem ha outra coisa que não é materia, e que é alguma coisa immaterial, permanente e independente da materia. Esta coisa é a \textsc{alma} ».}.

%%>>> t49

\pg{48} \tese{42}{Consequencias do materialismo.}
O \emph{materialismo} não só é falso nas doutrinas, mas é pernicioso nas consequencias; pois destroe a \emph{sciencia}, a \emph{arte}, a \emph{virtude}, a \emph{religião}, e a \emph{sociedade}.

a) Destroe a \emph{sciencia}. A sciencia é essencialmente dotada de \emph{unidade} e \emph{universalidade}. A \emph{unidade} é um effeito da \emph{generalisação; a universalidade,} da \emph{abstracção}. Ora a intelligencia, se fosse \pg{50} uma faculdade material ou organica, como ensina o materialismo, não poderia \emph{generalisar} nem \emph{abstrahir}; pois a operação não pode exceder a virtude do principio de que emana.

b) Destroe a \emph{arte}. A \emph{arte} não póde existir nem progredir, se o artista não formar um \emph{ideal}, pelo qual transforme a materia. Ora, o artista não poderia formar o \emph{ideal}, se não possuisse \pg{52} uma alma \emph{espiritual}, ou independente da materia; pois o \emph{ideal} forma-se na contemplação de ideas \emph{universaes}, que só podem ser concebidas por um principio espiritual.

c) Destroe a \emph{virtude}. O homem não poderia ser ornado de virtude, se não fosse \emph{livre} em praticar o bem e o mal. Ora o \emph{materialismo}, não reconhecendo no homem a liberdade \pg{54} de praticar o bem e o mal, nega a differença entre a virtude e o vicio, e assim abre o caminho a todos os crimes.

d) Destroe a \emph{Religião}. As verdades fundamentaes da Religião \pg{55} são a existencia de \textsc{Deus} e a immortalidade da alma humana. Ora o \emph{materialismo} nega essas duas verdades. --- Nega a existencia de \textsc{Deus}, porque só admitte o que está ao alcance \pg{56} dos sentidos, e \textsc{Deus} não se vê com os olhos corporeos. --- Nega a \emph{immortalidade} da alma humana; porque, se a alma fosse \pg{57} uma funcção do cerebro ou uma qualidade da materia, na dissolução do corpo, tambem a alma deveria acabar.

e) Destroe a \emph{sociedade}. A sociedade não póde existir nem alcançar o seu fim, se não se basear nos sentimentos da justiça, da moralidade e do bem. Ora, na doutrina do \emph{materialismo}, \pg{59}  que nega a existencia de \textsc{Deus}, a espiritualidade, a liberdade e a immortalidade da alma humana, é impossivel encontrar a base da justiça, da moralidade e do bem \footnote{Os materialistas, não podendo negar as consequencias do seu systema, acceitam o materialismo theorico e rejeitam o pratico.\par O procedimento dos adversarios é illogico. A sciencia estabelece os seus principios e as suas formulas, que em seguida a pratica traduz nas operações ; e por isso, quem admittir os principios, não póde deixar de admittir as consequencias. --- Assim como o materialismo pratico provoca o theorico, assim tambem este se reflecte no pratico, promovendo, sob um apparato scientifico, a desordem das paixões na vida individual e social.\par Concluiremos a critica do \emph{materialismo} com as palavras de \emph{Margerie}: « Emquanto a mim, nunca senti a minha alma tão serena na posse da verdade, como ao acabar este estudo sobre uma doutrina (\emph{materialista}), que se apresenta com uma confiança tão altiva em nome do espirito moderno, e que, todavia, abstrahindo até das suas consequencias, fortifica o espiritualismo pela impotencia dos seus ataques contra este. Na essencia, o materialismo de Büchner não é mais temivel que o de Broussais, de Cabanis, de Lamettrie, de Hobbes, de Epicuro ; e seria um grande signal de decadencia intellectual, de que \textsc{Deus} nos livre, se elle fizesse parte da philosophia. O seu logar é na sophistica; ahi ficará » (\emph{Phil}. \emph{contemp}., pag. 216).}.

\pg{49} \pg{51} \pg{53} \pg{58}
